% ============================================================
% ATUM_ATOM CODEX — arXiv submission, top-level file
% ============================================================
% Build:  make pdf          (from repository root)
% Pack:   ./arxiv-prep.sh   (produces build/arxiv/*.tar.gz)
%
% This file is pdflatex-safe. See src/preamble.tex for why.
% ============================================================

\documentclass[11pt,a4paper]{article}

% ============================================================
% ATUM_ATOM — Shared preamble
% ============================================================
% NOTE: This preamble is deliberately pdflatex-safe. The codex
% layout files under codex/layout/ use fontspec and therefore
% require XeLaTeX or LuaLaTeX; the arXiv submission in src/ does
% not, so that it builds under arXiv's default pdflatex path.
% Do not add fontspec here.
% ============================================================

\usepackage[T1]{fontenc}
\usepackage[utf8]{inputenc}
\usepackage{lmodern}
\usepackage{microtype}

\usepackage{amsmath}
\usepackage{amssymb}
\usepackage{amsthm}

\usepackage{graphicx}
\usepackage{booktabs}
\usepackage{array}

\usepackage{listings}
\usepackage{xcolor}
\usepackage{enumitem}

\usepackage[hidelinks]{hyperref}
\usepackage{cleveref}

% --- Theorem-like environments ---
\theoremstyle{definition}
\newtheorem{definition}{Definition}[section]
\newtheorem{protocol}[definition]{Protocol}

\theoremstyle{plain}
\newtheorem{property}[definition]{Property}

\theoremstyle{remark}
\newtheorem{remark}[definition]{Remark}

% --- Code listings ---
\lstdefinestyle{shell}{
  basicstyle=\ttfamily\footnotesize,
  breaklines=true,
  breakatwhitespace=false,
  columns=fullflexible,
  frame=single,
  framesep=4pt,
  rulecolor=\color{black!35},
  backgroundcolor=\color{black!3},
  showstringspaces=false,
  captionpos=b,
}
\lstset{style=shell}

% --- Semantic macros for the three pillars ---
% Used throughout so the vocabulary stays consistent with the
% codex (codex/book/codex_master.tex, codex/seals/).
\newcommand{\Manifest}{\textsc{Manifest}}
\newcommand{\Signature}{\textsc{Signature}}
\newcommand{\Seal}{\textsc{Seal}}
\newcommand{\MetaSeal}{\textsc{Meta-Seal}}
\newcommand{\Operator}{\textsc{Operator}}

% Hash of an object; signature of an object under key k.
\newcommand{\hash}[1]{\mathrm{H}(#1)}
\newcommand{\sign}[2]{\mathrm{Sign}_{#2}(#1)}

% Long hex digest, set small so wide tables stay in the margin.
\newcommand{\digest}[1]{{\ttfamily\scriptsize #1}}

% --- Document metadata (single source of truth) ---
\newcommand{\bindingcycle}{2026-06-21}
\newcommand{\atumoperator}{James Paul Stanley Jr}
\newcommand{\operatorfingerprint}{F9B445CFDBA1ECD6ECC1899CB58390F86B642281}
\newcommand{\codexroot}{\texttt{Hydrogenesi/OD\_Dragon\_Consciousness}}


\title{The Triadic Seal: A Three-Pillar Protocol for\\
       Structure, Attestation, and State in\\
       Multi-Repository Governance}

\author{\atumoperator\\
        \texttt{infinitysend@outlook.com}\\
        \small Hydrogenesi / ATUM\_ATOM Codex Root}

\date{Binding Cycle \bindingcycle}

\begin{document}

\maketitle

\begin{abstract}
Distributed research projects routinely span many repositories with no
single artifact that records what the collection was supposed to contain,
who vouched for it, or what state it was actually in at a given moment.
We describe the \emph{Triadic Seal}, a lightweight governance protocol that
binds three separable concerns --- structure, integrity, and state --- into a
single verifiable record. The \Manifest{} pillar fixes the authoritative
inventory of the collection under a content hash; the \Signature{} pillar
attests to that inventory with a public-key signature from a named operator;
the \Seal{} pillar snapshots observed repository state and reports drift
against the previous binding. A \MetaSeal{} consolidates the three into one
timestamped record committed to an append-only history. We give the
construction, the monthly binding cycle, the drift predicates, the threat
model the protocol does and does not address, and a reference implementation
in the ATUM\_ATOM Codex. The protocol is deliberately modest: it provides
detection and attribution, not prevention, and its guarantees rest entirely
on the secrecy of one operator key and the integrity of the commit history.
\end{abstract}

\tableofcontents

% ============================================================
\section{Introduction}
% ============================================================

A research program that outgrows a single repository acquires a problem it
did not have before: there is no longer any artifact that says what the
program consists of. The set of repositories becomes tacit knowledge. When a
repository is renamed, abandoned, forked, or quietly modified, nothing in the
system notices, because nothing in the system was ever told what the correct
configuration was.

The usual responses are either too heavy or too weak. Monorepos solve the
problem by eliminating it, at the cost of forcing unrelated work into one
history. Package manifests and lockfiles solve it for dependency graphs, but
they describe code that is consumed, not repositories that are governed.
Signed tags attest to a point in one repository's history and say nothing
about the collection.

The Triadic Seal takes a narrower aim. It answers three questions, and keeps
the answers separate so that each can fail independently and visibly:

\begin{enumerate}[label=(\roman*),leftmargin=2.2em]
  \item \textbf{What is supposed to be here?} --- the \Manifest{}.
  \item \textbf{Who says so?} --- the \Signature{}.
  \item \textbf{What is actually here, and how far has it moved?} --- the
        \Seal{}.
\end{enumerate}

Keeping these apart is the whole design. A conflated ``signed snapshot''
cannot distinguish a legitimate change that was never recorded from a
recorded change that was never authorised; the triadic split makes both
conditions nameable and separately detectable.

\paragraph{Contributions.}
\Cref{sec:related} places the protocol among existing supply-chain and
transparency systems. \Cref{sec:pillars} defines the three pillars and their
hash and signature relationships. \Cref{sec:metaseal} gives the \MetaSeal{}
consolidation.
\Cref{sec:cycle} specifies the binding cycle as an operational procedure.
\Cref{sec:drift} states the drift predicates. \Cref{sec:threat} sets out the
threat model, including the several things the protocol explicitly does not
defend against. \Cref{sec:impl} describes the reference implementation.
\Cref{sec:limitations} is candid about the protocol's limits.

% ============================================================
\section{Related Work}
\label{sec:related}
% ============================================================

The Triadic Seal borrows freely and claims little novelty in its primitives.
Its debts, and the boundaries between it and better-developed systems, are
worth stating precisely.

\paragraph{Hash-linked histories and timestamping.}
Committing to a document by publishing its digest, and chaining such
commitments over time, dates to Haber and Stornetta~\cite{haber1991};
Merkle's tree construction~\cite{merkle1988} underlies the efficient variants.
The meta-seal is a degenerate case: a linear sequence of commitments with no
tree and no external witness. Certificate Transparency~\cite{rfc6962}
demonstrates what the mature form looks like --- an append-only log with
verifiable inclusion and consistency proofs --- and the gap between that and a
sequence of commits in one repository is exactly the gap identified as N2 in
\cref{sec:nonthreats}.

\paragraph{Software supply-chain frameworks.}
TUF~\cite{tuf2010} addresses a strictly harder problem: it assumes key
compromise will happen and structures roles and thresholds so that it is
survivable. The Triadic Seal has no such structure and does not survive
operator key compromise at all. in-toto~\cite{intoto2019} attests to the
steps of a supply chain and who performed them, which is closer in spirit to
the \Signature{} pillar, but in-toto governs a \emph{build process} whereas
the Triadic Seal governs \emph{collection membership} --- a coarser and much
less demanding target. Sigstore~\cite{sigstore2022} removes long-lived key
custody from the signer through ephemeral certificates and a transparency
log; adopting it would substantially reduce the risk described in N1.

\paragraph{Positioning.}
None of these systems answers the specific question the Triadic Seal is built
for --- ``what repositories is this research program made of, and has that set
changed without anyone saying so?'' --- but each is stronger than the Triadic
Seal in the areas where they overlap. The protocol's justification is
operational cost, not cryptographic sophistication: it is small enough that a
single researcher will actually run it every month.

% ============================================================
\section{The Three Pillars}
\label{sec:pillars}
% ============================================================

Let $R = \{r_1,\dots,r_n\}$ be the governed collection of repositories, and
let $\mathrm{H}$ denote SHA-256, as specified in FIPS 180-4~\cite{fips180-4}.

\subsection{Pillar I --- \Manifest{} (Structure)}

\begin{definition}[Manifest]
\label{def:manifest}
A \emph{manifest} $M$ is a canonically serialised record listing, for each
$r_i \in R$: its name, its declared role, its governance boundary, and its
integration status. The \emph{manifest hash} is $\hash{M}$.
\end{definition}

The manifest is normative, not descriptive. It states what the collection
ought to be. Its hash is the fixed-point baseline against which every
subsequent observation is compared, so canonical serialisation matters: field
order, line endings, and encoding must be fixed, or $\hash{M}$ will change
for reasons that have nothing to do with governance.

\begin{remark}
The reference implementation additionally emits a \emph{fixity manifest} --- a
per-file SHA-256 listing of the Codex Root working tree, hashed as a whole.
This is a different object from $M$: the fixity manifest describes one
repository's bytes, whereas $M$ describes the collection's membership. Both
are useful; conflating them is a common and costly error, because a fixity
manifest changes on every commit while $M$ should change only when membership
or roles change.
\end{remark}

\subsection{Pillar II --- \Signature{} (Integrity)}

\begin{definition}[Signature]
\label{def:signature}
The \emph{signature} is $\sigma = \sign{M}{k}$, where $k$ is the private half
of the \Operator{}'s long-lived signing key. Verification uses the
corresponding public key and succeeds only if $M$ is bit-identical to the
signed input.
\end{definition}

The signature supplies two things the hash alone cannot: \emph{attribution}
(a named party asserted this configuration) and \emph{intent} (the assertion
was deliberate, not a side effect of a script run). A hash proves that two
byte strings match; only a signature proves that somebody meant it.

The reference deployment uses an Ed25519 key~\cite{rfc8032} held by the
\Operator{} and identified in published records by its OpenPGP fingerprint
\digest{\operatorfingerprint}.

\begin{remark}[On identifiers]
\label{rem:identifier}
A fingerprint is not a public key. The 160-bit value above identifies the key
but cannot verify anything on its own; verification requires the 32-byte
Ed25519 public key itself, distributed separately (\cref{sec:impl}). Records
that print a fingerprint next to the word ``key'' invite the reader to
believe they hold verification material when they do not. We keep the two
terms distinct throughout.
\end{remark}

\subsection{Pillar III --- \Seal{} (State)}

\begin{definition}[Seal]
\label{def:seal}
A \emph{seal} $S_t$ at binding time $t$ is a record of \emph{observed} state:
for each $r_i$, its last-activity timestamp and commit count since the
previous binding; plus the aggregate drift report of \cref{sec:drift} and the
\Operator{}'s attestation. Its digest is $\hash{S_t}$.
\end{definition}

Where the manifest is normative, the seal is empirical. The seal is the only
pillar that is expected to differ from one cycle to the next; a seal that
never changes indicates a collection that is either finished or unobserved.

% ============================================================
\section{The Triadic Meta-Seal}
\label{sec:metaseal}
% ============================================================

\begin{definition}[Meta-Seal]
The \emph{meta-seal} at time $t$ is
\[
  T_t \;=\; \bigl\langle\, t,\; \hash{M},\; \hash{\sigma},\; \hash{S_t},\;
              \mathrm{status},\; \mathrm{attestation} \,\bigr\rangle ,
\]
where $\mathrm{status} \in \{\textsc{triadic-stable},
\textsc{drift-detected}, \textsc{compromised}\}$.
\end{definition}

$T_t$ records the three pillar digests rather than the pillars themselves, so
it stays small and readable while remaining a complete commitment: any change
to $M$, $\sigma$, or $S_t$ changes $T_t$.

\begin{property}[Binding]
If $T_t$ verifies and the underlying $M$, $\sigma$, $S_t$ hash to the
recorded values, then at time $t$ the \Operator{} asserted structure $M$ and
observed state $S_t$ together. Neither can be substituted without detection.
\end{property}

\begin{property}[Non-prevention]
\label{prop:nonprevention}
$T_t$ constrains nothing about $t' > t$. An adversary with write access to a
governed repository can make arbitrary changes after a binding; the protocol
guarantees only that the \emph{next} binding will observe them. The Triadic
Seal is a detection mechanism, not an access control mechanism.
\end{property}

\Cref{prop:nonprevention} is stated as a property rather than buried in
limitations because it is the single most likely thing for a reader to get
wrong.

\begin{table}[t]
\centering
\caption{Meta-seal $T_{t_0}$ for the initial binding, cycle \bindingcycle.}
\label{tab:binding}
\begin{tabular}{@{}lll@{}}
\toprule
\textbf{Pillar} & \textbf{Digest (SHA-256)} & \textbf{Status} \\
\midrule
I --- \Manifest{}  & \digest{E6C46688C9883703B4F2B324DC3248A4} & VERIFIED \\
                   & \digest{E77E6D751746DCD87A509D1DBB2762D9} &          \\
\addlinespace
II --- \Signature{} & \digest{66FAB65A481E0FE3687CC839B7A33141} & VERIFIED \\
                   & \digest{1D8401DA69BA18EE42AA4513B7CA5751} &          \\
\addlinespace
III --- \Seal{}     & \digest{58B16B0280F3DF2B4A463A1647AE606E} & VERIFIED \\
                   & \digest{85B11E460FDD7092717A5D4BDF9145FD} &          \\
\bottomrule
\end{tabular}
\end{table}

\Cref{tab:binding} reproduces the initial binding recorded in the Codex Root.
Digests are split across two lines for typesetting only; the recorded values
are the 64-character concatenations.

% ============================================================
\section{The Binding Cycle}
\label{sec:cycle}
% ============================================================

Bindings are periodic, not event-driven. A monthly cadence is used in the
reference deployment. The trade-off is explicit: a longer period means a
longer window in which an undetected change can persist, while a shorter
period means more operator signing events and more opportunities for the
signing key to be exposed on an online machine.

\begin{protocol}[Monthly binding]
\label{prot:binding}
\begin{enumerate}[leftmargin=2.4em]
  \item \textbf{Collect.} For each $r_i \in R$, record last-activity
        timestamp and commit count since the previous binding.
  \item \textbf{Update.} Regenerate $M$ from current membership. If $M$
        changed, the change must be independently justified before signing;
        an unexplained manifest change is itself a drift indicator
        (\cref{sec:drift}).
  \item \textbf{Generate.} Emit $S_t$ with the drift report.
  \item \textbf{Sign.} Produce $\sigma = \sign{M}{k}$ with the operator key.
  \item \textbf{Consolidate.} Emit $T_t$ and commit all four artifacts to the
        Codex Root in a single commit.
\end{enumerate}
\end{protocol}

Step 5 is what makes the history append-only in practice rather than in
principle: the artifacts are only as immutable as the commit history that
carries them, which is why the Codex Root's history is the protocol's real
trust anchor (\cref{sec:threat}).

% ============================================================
\section{Drift Detection}
\label{sec:drift}
% ============================================================

\emph{Drift} is divergence between the normative manifest and observed
reality. The following predicates each raise a drift flag on repository
$r_i$ at binding time $t$:

\begin{enumerate}[label=D\arabic*.,leftmargin=2.6em]
  \item \textbf{Excess activity.} Commit count since the previous binding
        exceeds a threshold $\theta$ ($\theta = 10$ in the reference
        deployment).
  \item \textbf{Dormancy.} No activity for more than $\delta$ days
        ($\delta = 30$).
  \item \textbf{Unverified structural change.} $\hash{M}$ changed without a
        corresponding justified manifest update in \cref{prot:binding}
        step~2.
  \item \textbf{Signature failure.} $\sigma$ does not verify against $M$
        under the operator public key.
  \item \textbf{State anomaly.} Observed repository state is inconsistent
        with its declared role or integration status.
\end{enumerate}

\begin{remark}[On D1 and D2]
D1 and D2 are heuristics about \emph{attention}, not about correctness. A
healthy, busy repository trips D1 every cycle; a finished one trips D2
forever. They are deliberately noisy tripwires whose purpose is to force the
operator to look, and they should be tuned per repository rather than
globally. D3 and D4 are the predicates with security content: D4 in
particular cannot be tripped by benign activity.
\end{remark}

A binding in which any predicate fires is recorded with status
\textsc{drift-detected} rather than \textsc{triadic-stable}. Drift does not
halt the cycle --- the binding still completes and is still signed --- because a
signed record of a known-drifted state is more useful than no record.

% ============================================================
\section{Threat Model}
\label{sec:threat}
% ============================================================

\subsection{What the protocol addresses}

\begin{enumerate}[label=T\arabic*.,leftmargin=2.6em]
  \item \textbf{Silent membership change.} A repository added to or removed
        from the collection without record. Detected at the next binding via
        $\hash{M}$ mismatch.
  \item \textbf{Retroactive tampering with governance records.} Modifying a
        past manifest or seal. Detected because $T_t$ commits to their
        digests and $T_t$ itself is in the commit history.
  \item \textbf{Repudiation.} An operator denying that a configuration was
        authorised. Prevented by $\sigma$, to the extent that $k$ is
        genuinely under sole operator control.
  \item \textbf{Unattended decay.} Repositories drifting out of the program
        without anyone noticing. Detected by D1--D2.
\end{enumerate}

\subsection{What the protocol does \emph{not} address}
\label{sec:nonthreats}

\begin{enumerate}[label=N\arabic*.,leftmargin=2.6em]
  \item \textbf{Operator key compromise.} An adversary holding $k$ can
        produce bindings indistinguishable from genuine ones. The protocol
        has no threshold or multi-party signing, so this is a total
        compromise, not a degraded mode. This is the dominant risk.
  \item \textbf{Compromise of the Codex Root history.} An adversary who can
        force-push to the Codex Root can rewrite the sequence of meta-seals.
        Independent anchoring (publishing $\hash{T_t}$ to a witness outside
        the repository) would mitigate this and is not currently done.
  \item \textbf{Between-binding changes.} By \cref{prop:nonprevention}, the
        entire inter-binding window is unprotected.
  \item \textbf{Content correctness.} The protocol says nothing about
        whether the code, data, or claims inside a governed repository are
        correct. A perfectly triadic-stable collection can be entirely
        wrong.
  \item \textbf{Observation integrity.} The seal records what the collection
        tool observed. An adversary who controls the observation environment
        controls $S_t$.
\end{enumerate}

N1 and N2 together mean the protocol's security reduces to two conventional
assumptions --- key secrecy and hosting-platform integrity --- neither of which
the protocol itself strengthens. Its contribution is organisational: it makes
the state of a sprawling collection legible and attributable at regular
intervals.

% ============================================================
\section{Reference Implementation}
\label{sec:impl}
% ============================================================

The ATUM\_ATOM Codex Root (\codexroot) implements the protocol with a small
set of scripts and a fixed directory layout:

\begin{itemize}[leftmargin=1.6em]
  \item \texttt{.triadic/manifest.txt} --- the manifest $M$.
  \item \texttt{.triadic/governance.md} --- the governance protocol, normative.
  \item \texttt{codex/seals/seal\_*.txt} --- seals $S_t$.
  \item \texttt{codex/seals/triadic\_seal\_*.txt} --- meta-seals $T_t$.
  \item \texttt{manifests/fixity\_manifest\_*.txt} --- per-file fixity records.
  \item \texttt{scripts/} --- manifest generation, binding, and export.
\end{itemize}

\subsection{Verification}

Verifying a published binding requires the operator's Ed25519 public key,
distributed independently of this repository. Given that key, a third party
verifies a binding as follows.

\noindent Confirm the meta-seal digest:
\begin{lstlisting}
sha256sum codex/seals/triadic_seal_20260621_104332.txt
\end{lstlisting}

\noindent Verify the manifest signature. For an OpenPGP-held Ed25519 key:
\begin{lstlisting}
gpg --verify .triadic/manifest.txt.sig .triadic/manifest.txt
\end{lstlisting}

\noindent For a raw Ed25519 key in PEM form:
\begin{lstlisting}
openssl pkeyutl -verify -pubin -inkey ed25519_pub.pem -rawin \
    -sigfile .triadic/manifest.txt.sig -in .triadic/manifest.txt
\end{lstlisting}

\begin{remark}[Verification command hygiene]
Ed25519 is a pure signature scheme: it hashes the message internally and must
be handed the message itself, not a digest. The command
\texttt{openssl dgst -sha256 -verify \dots}, which appears in some earlier
ATUM\_ATOM documentation, describes a hash-then-sign workflow appropriate to
RSA or ECDSA and will not verify an Ed25519 signature. The forms above are
the correct ones. Publishing a verification command that cannot succeed is
worse than publishing none, because a reader who runs it and sees a failure
learns nothing about whether the artifact is genuine.
\end{remark}

\subsection{Governed collection}

\begin{table}[t]
\centering
\caption{Governed repositories at cycle \bindingcycle.}
\label{tab:repos}
\begin{tabular}{@{}lll@{}}
\toprule
\textbf{Repository} & \textbf{Declared role} & \textbf{Status} \\
\midrule
Phoenix\_Ignition\_TOE      & Ceremonial framework, visualisation & ACTIVE \\
Phoenix-2.0-Apex-Edition    & Quantum modelling engine            & ACTIVE \\
UNI\_VERSE-                 & Universe-scale architecture         & ACTIVE \\
TUA                         & Unified theoretical framework       & ACTIVE \\
TOE                         & Theory-of-everything research       & ACTIVE \\
Quantum                     & Quantum computing research          & ACTIVE \\
\bottomrule
\end{tabular}
\end{table}

\Cref{tab:repos} lists the six governed repositories. ``Declared role'' is
the operator's assertion recorded in $M$; it is not independently validated
by the protocol, and \cref{sec:nonthreats} N4 applies.

% ============================================================
\section{Limitations and Future Work}
\label{sec:limitations}
% ============================================================

\paragraph{Single point of failure.}
One key, one operator, no threshold scheme. Splitting signing authority
across $m$-of-$n$ holders, along the lines of TUF's role model~\cite{tuf2010},
would remove the dominant risk of N1 at the cost of a materially heavier
binding ceremony. Adopting keyless signing~\cite{sigstore2022} is the
lighter-weight alternative.

\paragraph{No external anchor.}
Meta-seal digests are recorded only inside the repository whose history they
are meant to protect. Publishing $\hash{T_t}$ to any independent, append-only
witness~\cite{rfc6962} would close N2. This is the cheapest available
improvement and should be done first.

\paragraph{Unvalidated observation.}
The collection step trusts its own execution environment (N5). Running it in
a reproducible environment and signing the observation transcript would
narrow the gap.

\paragraph{Coarse drift predicates.}
D1 and D2 use global thresholds where per-repository baselines would be more
informative. Repositories differ enormously in expected commit rate, and a
threshold that is a useful tripwire for one is pure noise for another.

\paragraph{No formal treatment.}
This paper gives an operational specification, not a proof. The binding
property of \cref{sec:metaseal} follows straightforwardly from collision
resistance and signature unforgeability, but is not stated as a reduction.

% ============================================================
\section{Conclusion}
% ============================================================

The Triadic Seal separates three questions that governance systems usually
answer together and therefore answer badly: what a collection should contain,
who vouched for it, and what it actually contains. Keeping the answers in
distinct, individually hashed pillars, consolidated into one signed and
committed meta-seal per cycle, yields a record that is cheap to produce and
readable by a third party years later.

The protocol's guarantees are narrow and worth restating plainly: it detects,
it does not prevent; it attributes, it does not validate; and it inherits
rather than improves the security of the key and the hosting platform beneath
it. Within those bounds it turns an undocumented collection of repositories
into something with a history that can be audited.

% ============================================================
\bibliographystyle{plain}
\bibliography{references}
% ============================================================

\end{document}
